\subsection{Forwarding and concurrency pilot}\label{sec:live-study}
A separate pilot fixes initial messages and compares \emph{Full} tool feedback with \emph{Projected} task-required fields. Projection removes supplementary records but keeps the customer request, checks and order fields, so some model-link recognition is expected by design. Category-bearing synthetic records and two fixed 8B models yield 3,120 local LangGraph/MCP executions over streaming TLS~1.3/HTTP/2. Each test cell has 120 tasks in twenty nonempty and four empty sessions. First-feedback LR changes from 43 to 23 correct for Qwen and 40 to 22 for Llama; none of the eight primary local tests rejects after Holm correction (minimum adjusted $p=.0508$). Complete projected model-view scores remain 44/120 and 33/120 (chance 20/120), consistent with those retained fields. Projection also changes length and execution order, so the pilot neither isolates feedback content nor, by itself, establishes Finding~1; a category-free projection is the needed control (Table~\ref{tab:open-questions}, Q2). Supplement S1 retains all phases, attacks, outcomes and uncertainty.

Concurrency changes the evaluation unit. Supplied windows mix simultaneously active tasks: joint-view recognition is already low with dedicated connections (30/120 and 20/120), before sharing (23 and 19) or joint shaping (20 and 22). Event-level tests remain informative. At IoU .5, learned-detector F1 is .117--.274 with joint shaping versus .077--.249 without, at 19.5--22.5$\times$ model bytes. Shared-condition precision is 4.3--18.6\%, with 8.21--44.79 false events per test session. A constant all-active baseline outranks learned detectors at IoU .3 and falls below them at .5. These diagnostics motivate reporting inference together with event discovery and matching-threshold sensitivity; they supply no general shaping or localization guarantee.
